% GENERATED FILE. Edit canonical lessons, then run npm run generate:books.
% canonical-source-hash: fnv1a64:778e9c18c9a93eed
% canonical-lessons: AR-C282-hasiba, AR-C282-miqlama, AR-C282-safha, AR-C282-jumla, AR-C282-hallaq, AR-R282-first-pass-words-from-chapters-278-279, AR-R282-second-pass-words-from-chapters-280-282

\chapter{Calculator, Pencil case, Page, Sentence, Barber}
\label{ch:ar-calculator-pencil-case-page-sentence-barber}
\hlchaptermodality{282}

\begin{chapteropening}
\textbf{By the end of this chapter:} \emph{I can say a calculator, a pencil case, a page, a sentence and a barber.}

Complete the last lesson of chapter 282: I can say a calculator, a pencil case, a page, a sentence and a barber.
\end{chapteropening}

\section[\texorpdfstring{\ar{حاسبة} (ḥāsiba)}{حاسبة (ḥāsiba)}]{\ar{حاسبة} (ḥāsiba) — a calculator}

\label{lesson:AR-C282-hasiba}



\begin{quote}
Before the new one: say the Arabic for a presenter, then the Arabic for an actor.
\end{quote}

\subsection*{You'll want to know: \ar{حاسبة}}
\textbf{\ar{حاسبة}} — \emph{ḥāsiba} — \textquotedblleft{}a calculator\textquotedblright{}.

\begin{cousinweb}
\textbf{\ar{حاسبة}} (\emph{ḥāsiba}) is built on the native Arabic root \emph{ḥ-s-b}, \textquotedblleft{}to count, to reckon\textquotedblright{} (Lane). \emph{Ḥāsiba} is the feminine active participle, \textquotedblleft{}a thing that counts\textquotedblright{}.
\end{cousinweb}

\begin{grammarlens}[title={What you've built}]
One more word for this chapter.
\end{grammarlens}

\subsection*{Your turn}
\begin{itemize}
\raggedright
  \item \emph{Say it:} \emph{ḥāsiba}
  \item \emph{Say it:} \emph{ḥāsiba}, once more
  \item \emph{Say it:} say \emph{ḥāsiba}
  \item \emph{Recall:} say the Arabic for a presenter, then the Arabic for an actor, then say \emph{ḥāsiba} again
\end{itemize}

\begin{tcolorbox}[breakable,colback=teal!4,colframe=teal!35!black,title={Before you move on}]
What does \ar{حاسبة} mean? (\textquotedblleft{}A calculator\textquotedblright{}.) Say it once more.
\end{tcolorbox}
\section[\texorpdfstring{\ar{مقلمة} (miqlama)}{مقلمة (miqlama)}]{\ar{مقلمة} (miqlama) — a pencil case}

\label{lesson:AR-C282-miqlama}



\begin{quote}
Before the new one: say the Arabic for an actor, then the Arabic for a calculator.
\end{quote}

\subsection*{You'll want to know: \ar{مقلمة}}
\textbf{\ar{مقلمة}} — \emph{miqlama} — \textquotedblleft{}a pencil case\textquotedblright{}.

\begin{cousinweb}
\textbf{\ar{مقلمة}} (\emph{miqlama}) is built on the native Arabic root \emph{q-l-m}, \textquotedblleft{}a reed pen\textquotedblright{} (Lane). \emph{Miqlama} is a case for pens.
\end{cousinweb}

\begin{grammarlens}[title={What you've built}]
One more word for this chapter.
\end{grammarlens}

\subsection*{Your turn}
\begin{itemize}
\raggedright
  \item \emph{Say it:} \emph{miqlama}
  \item \emph{Say it:} \emph{miqlama}, once more
  \item \emph{Say it:} say \emph{miqlama}
  \item \emph{Recall:} say the Arabic for an actor, then the Arabic for a calculator, then say \emph{miqlama} again
\end{itemize}

\begin{tcolorbox}[breakable,colback=teal!4,colframe=teal!35!black,title={Before you move on}]
What does \ar{مقلمة} mean? (\textquotedblleft{}A pencil case\textquotedblright{}.) Say it once more.
\end{tcolorbox}
\section[\texorpdfstring{\ar{صفحة} (ṣafḥa)}{صفحة (ṣafḥa)}]{\ar{صفحة} (ṣafḥa) — a page}

\label{lesson:AR-C282-safha}



\begin{quote}
Before the new one: say the Arabic for a calculator, then the Arabic for a pencil case.
\end{quote}

\subsection*{You'll want to know: \ar{صفحة}}
\textbf{\ar{صفحة}} — \emph{ṣafḥa} — \textquotedblleft{}a page\textquotedblright{}.

\begin{cousinweb}
\textbf{\ar{صفحة}} (\emph{ṣafḥa}) is built on the native Arabic root \emph{ṣ-f-ḥ}, \textquotedblleft{}a broad surface; the face\textquotedblright{} (Lane). \emph{Ṣafḥa} is a broad face: a page.
\end{cousinweb}

\begin{grammarlens}[title={What you've built}]
One more word for this chapter.
\end{grammarlens}

\subsection*{Your turn}
\begin{itemize}
\raggedright
  \item \emph{Say it:} \emph{ṣafḥa}
  \item \emph{Say it:} \emph{ṣafḥa}, once more
  \item \emph{Say it:} say \emph{ṣafḥa}
  \item \emph{Recall:} say the Arabic for a calculator, then the Arabic for a pencil case, then say \emph{ṣafḥa} again
\end{itemize}

\begin{tcolorbox}[breakable,colback=teal!4,colframe=teal!35!black,title={Before you move on}]
What does \ar{صفحة} mean? (\textquotedblleft{}A page\textquotedblright{}.) Say it once more.
\end{tcolorbox}
\section[\texorpdfstring{\ar{جملة} (jumla)}{جملة (jumla)}]{\ar{جملة} (jumla) — a sentence}

\label{lesson:AR-C282-jumla}



\begin{quote}
Before the new one: say the Arabic for a pencil case, then the Arabic for a page.
\end{quote}

\subsection*{You'll want to know: \ar{جملة}}
\textbf{\ar{جملة}} — \emph{jumla} — \textquotedblleft{}a sentence\textquotedblright{}.

\begin{cousinweb}
\textbf{\ar{جملة}} (\emph{jumla}) is built on the native Arabic root \emph{j-m-l}, \textquotedblleft{}the whole, the sum\textquotedblright{} (Lane). \emph{Jumla} is a sum, a whole: a sentence.
\end{cousinweb}

\begin{grammarlens}[title={What you've built}]
One more word for this chapter.
\end{grammarlens}

\subsection*{Your turn}
\begin{itemize}
\raggedright
  \item \emph{Say it:} \emph{jumla}
  \item \emph{Say it:} \emph{jumla}, once more
  \item \emph{Say it:} say \emph{jumla}
  \item \emph{Recall:} say the Arabic for a pencil case, then the Arabic for a page, then say \emph{jumla} again
\end{itemize}

\begin{tcolorbox}[breakable,colback=teal!4,colframe=teal!35!black,title={Before you move on}]
What does \ar{جملة} mean? (\textquotedblleft{}A sentence\textquotedblright{}.) Say it once more.
\end{tcolorbox}
\section[\texorpdfstring{\ar{حلاق} (ḥallāq)}{حلاق (ḥallāq)}]{\ar{حلاق} (ḥallāq) — a barber}

\label{lesson:AR-C282-hallaq}



\begin{quote}
Before the new one: say the Arabic for a page, then the Arabic for a sentence.
\end{quote}

\subsection*{You'll want to know: \ar{حلاق}}
\textbf{\ar{حلاق}} — \emph{ḥallāq} — \textquotedblleft{}a barber\textquotedblright{}.

\begin{cousinweb}
\textbf{\ar{حلاق}} (\emph{ḥallāq}) is built on the native Arabic root \emph{ḥ-l-q}, \textquotedblleft{}to shave\textquotedblright{} (Lane). \emph{Ḥallāq} is the \emph{faʿʿāl} trade noun, \textquotedblleft{}barber\textquotedblright{}.
\end{cousinweb}

\begin{grammarlens}[title={What you've built}]
One more word for this chapter.
\end{grammarlens}

\subsection*{Your turn}
\begin{itemize}
\raggedright
  \item \emph{Say it:} \emph{ḥallāq}
  \item \emph{Say it:} \emph{ḥallāq}, once more
  \item \emph{Say it:} say \emph{ḥallāq}
  \item \emph{Recall:} say the Arabic for a page, then the Arabic for a sentence, then say \emph{ḥallāq} again
\end{itemize}

\begin{tcolorbox}[breakable,colback=teal!4,colframe=teal!35!black,title={Before you move on}]
What does \ar{حلاق} mean? (\textquotedblleft{}A barber\textquotedblright{}.) Say it once more.
\end{tcolorbox}
\section[(dialogue)]{Words from chapters 278-279 — a first pass}

\label{lesson:AR-R282-first-pass-words-from-chapters-278-279}



\begin{quote}
Say the words for a sentence and a barber.
\end{quote}

\subsection*{Your turn}
Say these aloud:

\begin{itemize}
\raggedright
  \item a kettle, a tray, a strainer, a ladle, a barrel
  \item a rectangle, a coat hanger, a sponge, a tape, glue
\end{itemize}

\begin{tcolorbox}[breakable,colback=teal!4,colframe=teal!35!black,title={Before you move on}]
A kettle? (\textbf{\ar{غلاية}}, \emph{ghallāya}.) A tray? (\textbf{\ar{صينية}}, \emph{ṣīniyya}.) A strainer? (\textbf{\ar{مصفاة}}, \emph{miṣfāh}.) A ladle? (\textbf{\ar{مغرفة}}, \emph{mighrafa}.) A barrel? (\textbf{\ar{برميل}}, \emph{barmīl}.) A rectangle? (\textbf{\ar{مستطيل}}, \emph{mustaṭīl}.) A coat hanger? (\textbf{\ar{مشجب}}, \emph{mishjab}.) A sponge? (\textbf{\ar{إسفنجة}}, \emph{isfanja}.) A tape? (\textbf{\ar{شريط}}, \emph{sharīṭ}.) Glue? (\textbf{\ar{غراء}}, \emph{ghirāʾ}.)
\end{tcolorbox}
\section[(dialogue)]{Words from chapters 280-282 — a second pass}

\label{lesson:AR-R282-second-pass-words-from-chapters-280-282}



\begin{quote}
Say the words for a client and an engineer.
\end{quote}

\subsection*{Your turn}
Say these aloud:

\begin{itemize}
\raggedright
  \item a client, an engineer, a lawyer, a pilot, a pharmacist
  \item a number, a phone call, information, a presenter, an actor
  \item a calculator, a pencil case, a page, a sentence, a barber
\end{itemize}

\begin{tcolorbox}[breakable,colback=teal!4,colframe=teal!35!black,title={Before you move on}]
A client? (\textbf{\ar{عميل}}, \emph{ʿamīl}.) An engineer? (\textbf{\ar{مهندس}}, \emph{muhandis}.) A lawyer? (\textbf{\ar{محامي}}, \emph{muḥāmī}.) A pilot? (\textbf{\ar{طيار}}, \emph{ṭayyār}.) A pharmacist? (\textbf{\ar{صيدلي}}, \emph{ṣaydalī}.) A number? (\textbf{\ar{رقم}}, \emph{raqm}.) A phone call? (\textbf{\ar{مكالمة}}, \emph{mukālama}.) Information? (\textbf{\ar{معلومات}}, \emph{maʿlūmāt}.) A presenter? (\textbf{\ar{مذيع}}, \emph{mudhīʿ}.) An actor? (\textbf{\ar{ممثل}}, \emph{mumaththil}.) A calculator? (\textbf{\ar{حاسبة}}, \emph{ḥāsiba}.) A pencil case? (\textbf{\ar{مقلمة}}, \emph{miqlama}.) A page? (\textbf{\ar{صفحة}}, \emph{ṣafḥa}.) A sentence? (\textbf{\ar{جملة}}, \emph{jumla}.) A barber? (\textbf{\ar{حلاق}}, \emph{ḥallāq}.)
\end{tcolorbox}
